%% evenend-sample-colsep.tex：evenendの見本（段組ごとの段間と段間罫）
%% マニュアル（evenend-ja.tex）が貼り込む
\documentclass[twocolumn,paper=a6,fontsize=8pt]{jlreq}
\usepackage{evenend}
\setlength\columnseprule{0.4pt}
\newcommand\para{吾輩は猫である．名前はまだ無い．どこで生れたかとんと見当がつかぬ．何でも薄暗いじめじめした所でニャーニャー泣いていた事だけは記憶している．\par}
\newcommand\short{吾輩は猫である．名前はまだ無い．どこで生れたかとんと見当がつかぬ．\par}
\begin{document}
\noindent\textgt{既定（段間2字，罫0.4pt）}\par
\short
\evenendcolumns[columnsep=1\zw,columnseprule=0pt]{3}
\noindent\textgt{段間1字・罫なし}\par
\para
\evenendcolumns[columnsep=4\zw,columnseprule=1.5pt]{2}
\noindent\textgt{段間4字・罫1.5pt}\par
\short
\evenendcolumns[columnsep=3\zw]{3}
\noindent\textgt{段間3字}\par
\para
\evenendcolumns{2}
\noindent\textgt{既定に戻る}\par
\para
\end{document}
